\documentclass[12pt,a4paper]{report}

% ── Packages ──────────────────────────────────────────────────────────────────
\usepackage[a4paper, left=1.25in, right=1in, top=1in, bottom=1in]{geometry}
\usepackage{times}
\usepackage{setspace}
\usepackage{graphicx}
\usepackage{booktabs}
\usepackage{array}
\usepackage{amsmath}
\usepackage{amssymb}
\usepackage{hyperref}
\usepackage{caption}
\usepackage{fancyhdr}
\usepackage{titlesec}
\usepackage{tocloft}
\usepackage{parskip}
\usepackage{enumitem}
\usepackage{xcolor}
\usepackage{microtype}
\usepackage{multirow}
\usepackage{longtable}
\usepackage{listings}
\usepackage{algorithm}
\usepackage{algpseudocode}
\usepackage{float}
\usepackage{tikz}

\hypersetup{colorlinks=true, linkcolor=black, citecolor=black, urlcolor=blue}

\onehalfspacing

% ── Chapter/Section formatting ─────────────────────────────────────────────────
\titleformat{\chapter}[display]
  {\normalfont\Large\bfseries\centering}
  {Chapter \thechapter}
  {10pt}
  {\large\centering}
\titlespacing*{\chapter}{0pt}{20pt}{30pt}

\titleformat{\section}{\normalfont\large\bfseries}{\thesection}{1em}{}
\titleformat{\subsection}{\normalfont\normalsize\bfseries}{\thesubsection}{1em}{}
\titleformat{\subsubsection}{\normalfont\normalsize\itshape}{\thesubsubsection}{1em}{}

% ── Header/Footer ─────────────────────────────────────────────────────────────
\pagestyle{fancy}
\fancyhf{}
\fancyhead[L]{\small\textit{ExplainHire --- B.Tech Major Project Report}}
\fancyhead[R]{\small\textit{DTU, 2025--26}}
\fancyfoot[C]{\thepage}
\renewcommand{\headrulewidth}{0.4pt}

% ── Listings style ────────────────────────────────────────────────────────────
\lstset{
  basicstyle=\small\ttfamily,
  breaklines=true,
  frame=single,
  backgroundcolor=\color{gray!10}
}

% ── Figure placeholder command ────────────────────────────────────────────────
\newcommand{\figplaceholder}[2][0.85\textwidth]{%
  \begin{tikzpicture}
    \draw[thick, dashed, gray] (0,0) rectangle (#1, 5cm);
    \node[gray] at (#1/2, 2.5cm) {\textit{[Figure: #2]}};
  \end{tikzpicture}%
}

% ══════════════════════════════════════════════════════════════════════════════
\begin{document}

% ── Cover Page ────────────────────────────────────────────────────────────────
\begin{titlepage}
\pagestyle{empty}
\begin{center}

\textbf{\large DEPT. OF COMPUTER SCIENCE \& ENGINEERING}\\[4pt]
\textbf{\large DELHI TECHNOLOGICAL UNIVERSITY}\\[2pt]
{\small (Formerly Delhi College of Engineering)}\\
{\small Bawana Road, Delhi--110042}\\[24pt]

\rule{\linewidth}{0.5pt}\\[12pt]

{\LARGE \textbf{ExplainHire: Neurosymbolic AI for}}\\[8pt]
{\LARGE \textbf{Explainable Resume--Job Matching}}\\[12pt]

\rule{\linewidth}{0.5pt}\\[24pt]

\textbf{\large A PROJECT REPORT}\\[6pt]
\textit{SUBMITTED IN PARTIAL FULFILLMENT OF THE REQUIREMENTS}\\[2pt]
\textit{FOR THE AWARD OF THE DEGREE}\\[6pt]
{\Large \textbf{OF}}\\[6pt]
{\Large \textbf{BACHELOR OF TECHNOLOGY}}\\[6pt]
{\Large \textbf{IN}}\\[6pt]
{\Large \textbf{COMPUTER SCIENCE \& ENGINEERING}}\\[30pt]

\textbf{\large Submitted by:}\\[10pt]
\begin{tabular}{c}
  {\large Yoosha Raza \hspace{4pt} (2K22/CO/521)} \\[6pt]
  {\large Yanshu \hspace{4pt} (2K22/CO/507)} \\[6pt]
  {\large Yogarth \hspace{4pt} (2K22/CO/519)} \\
\end{tabular}\\[28pt]

\textbf{\large Under the supervision of}\\[8pt]
{\large \textbf{Gull Kaur}}\\[4pt]
Assistant Professor\\
Department of Computer Science \& Engineering\\
Delhi Technological University\\[28pt]

\rule{\linewidth}{0.5pt}\\[8pt]
{\large \textbf{MAY 2026}}\\[4pt]
\rule{\linewidth}{0.5pt}

\end{center}
\end{titlepage}

% ── Candidate's Declaration ───────────────────────────────────────────────────
\newpage
\pagestyle{plain}
\pagenumbering{roman}
\setcounter{page}{1}
\thispagestyle{plain}

\begin{center}
\textbf{\large DEPT. OF COMPUTER SCIENCE \& ENGINEERING}\\
\textbf{\large DELHI TECHNOLOGICAL UNIVERSITY}\\
{\small (Formerly Delhi College of Engineering)}\\
{\small Bawana Road, Delhi--110042}\\[20pt]
\textbf{\Large CANDIDATE'S DECLARATION}\\[4pt]
\rule{0.5\linewidth}{0.4pt}
\end{center}

\vspace{14pt}

We, \textbf{Yoosha Raza (2K22/CO/521)}, \textbf{Yanshu (2K22/CO/507)}, and \textbf{Yogarth (2K22/CO/519)}, students of B.Tech (Computer Science \& Engineering), Batch 2022--2026, hereby declare that the Project Dissertation titled \textit{``ExplainHire: Neurosymbolic AI for Explainable Resume--Job Matching''}, submitted to the Department of Computer Science \& Engineering, Delhi Technological University, is an original work carried out by us. This work has not previously formed the basis for the award of any Degree, Diploma, or similar title by this or any other institution. All sources have been duly cited and referenced.

The implementation, experiments, results, and analysis presented in this report are the result of our own independent work, conducted between August 2025 and May 2026, under the supervision of \textbf{Gull Kaur}, Assistant Professor, CSE, DTU.

\vspace{24pt}
\noindent Place: New Delhi \hfill Date: May 2026

\vspace{36pt}
\begin{tabular}{p{4cm} p{4cm} p{4cm}}
\textbf{Yoosha Raza} & \textbf{Yanshu} & \textbf{Yogarth} \\
(2K22/CO/521) & (2K22/CO/507) & (2K22/CO/519) \\
\end{tabular}

\vspace*{\fill}
\begin{center}{\small \textit{i}}\end{center}

% ── Certificate ───────────────────────────────────────────────────────────────
\newpage
\thispagestyle{plain}

\begin{center}
\textbf{\large DEPT. OF COMPUTER SCIENCE \& ENGINEERING}\\
\textbf{\large DELHI TECHNOLOGICAL UNIVERSITY}\\
{\small (Formerly Delhi College of Engineering)}\\
{\small Bawana Road, Delhi--110042}\\[20pt]
\textbf{\Large CERTIFICATE}\\[4pt]
\rule{0.5\linewidth}{0.4pt}
\end{center}

\vspace{14pt}

This is to certify that the project report titled \textit{``ExplainHire: Neurosymbolic AI for Explainable Resume--Job Matching''}, submitted by \textbf{Yoosha Raza (2K22/CO/521)}, \textbf{Yanshu (2K22/CO/507)}, and \textbf{Yogarth (2K22/CO/519)} in partial fulfillment of the requirements for the award of the degree of \textbf{Bachelor of Technology in Computer Science \& Engineering} from Delhi Technological University, is a bonafide record of the project work carried out by them under my guidance during the academic year 2025--26.

The project work is original and has not been submitted to any other university or institution. I recommend this report for acceptance.

\vspace{30pt}
\noindent Place: New Delhi \hfill Date: May 2026

\vspace{36pt}
\noindent\textbf{Gull Kaur}\\
(Supervisor)\\
Assistant Professor, CSE, DTU, Delhi--110042

\vspace*{\fill}
\begin{center}{\small \textit{ii}}\end{center}

% ── Abstract ──────────────────────────────────────────────────────────────────
\newpage
\thispagestyle{plain}

\begin{center}
\textbf{\Large ABSTRACT}\\[4pt]
\rule{0.5\linewidth}{0.4pt}
\end{center}

\vspace{10pt}

\noindent\textbf{Keywords---} Neurosymbolic AI, Resume--Job Matching, Skill Ontology, Sentence-BERT, XGBoost, SHAP Explainability, Applicant Tracking System, NLP

\vspace{12pt}

Existing Applicant Tracking Systems (ATS) rely on keyword co-occurrence to screen resumes, causing two critical failures: \textit{semantic blindness} (equivalent skills described differently score zero) and \textit{decision opacity} (no structured justification for outcomes). This report presents \textbf{ExplainHire}, a neurosymbolic AI system that resolves both.

ExplainHire integrates three complementary matching signals in a seven-layer pipeline: (1) a \textbf{434-node symbolic skill ontology graph} derived from O*NET, enabling six-level weighted traversal capturing exact, alias, and ontological proximity matches; (2) a \textbf{domain-adapted SBERT model} (\texttt{all-MiniLM-L6-v2}) fine-tuned on 887 real resume--JD pairs from 144 genuine B.Tech student resumes matched against 17 industry job descriptions; and (3) \textbf{structural features} capturing years-of-experience gap, education level, and resume section coverage. The three signals are decomposed into 14 numeric features and fused by an XGBoost classifier trained on 1,854 labeled pairs.

Explainability is delivered through SHAP TreeExplainer (feature-level attribution) and a sentence-level evidence mapper (exact resume sentence proving each matched skill), complemented by a counterfactual skill gap ranker.

ExplainHire achieves a 5-fold stratified CV F1 of \textbf{0.9045} and a full-dataset F1 of \textbf{0.951}, representing \textbf{+11.86\%} improvement over the best single-signal baseline (Graph + SBERT, F1 = 0.832). On the public ResumeAtlas benchmark (Heakl et al., 2024), it achieves macro F1 of \textbf{0.720} vs. a reported TF-IDF baseline of 0.61 (\textbf{+11.06\%}). ExplainHire is, to our knowledge, the first resume--job matching system to provide structured, auditable, sentence-level explanations combined with ranked skill gap analysis.

\vspace*{\fill}
\begin{center}{\small \textit{iii}}\end{center}

% ── Acknowledgement ───────────────────────────────────────────────────────────
\newpage
\thispagestyle{plain}

\begin{center}
\textbf{\Large ACKNOWLEDGEMENT}\\[4pt]
\rule{0.5\linewidth}{0.4pt}
\end{center}

\vspace{12pt}

We are deeply grateful to our project supervisor, \textbf{Gull Kaur}, Assistant Professor, CSE, DTU, for her invaluable guidance throughout this project. Her insight into artificial intelligence and natural language processing shaped every stage of our work.

We thank the \textbf{Head of Department} and all faculty of the CSE Department, DTU, for creating a research environment that supports ambitious final-year projects. We owe a special debt to our \textbf{144 fellow B.Tech students} who generously contributed their resumes --- our results rest on their participation. We also thank the creators of O*NET, the sentence-transformers library, the ResumeAtlas benchmark, and the Kaggle resume dataset.

Finally, we thank our \textbf{families and parents}, whose unwavering support through four years of engineering education made everything possible.

\vspace{36pt}
\begin{tabular}{p{4cm} p{4cm} p{4cm}}
\textbf{Yoosha Raza} & \textbf{Yanshu} & \textbf{Yogarth} \\
(2K22/CO/521) & (2K22/CO/507) & (2K22/CO/519) \\
\end{tabular}

\vspace*{\fill}
\begin{center}{\small \textit{iv}}\end{center}

% ── Table of Contents ─────────────────────────────────────────────────────────
\newpage
\tableofcontents

\newpage
\listoffigures

\newpage
\listoftables

% ── Abbreviations ─────────────────────────────────────────────────────────────
\newpage
\begin{center}
\textbf{\Large LIST OF SYMBOLS, ABBREVIATIONS AND NOMENCLATURE}\\[4pt]
\rule{0.5\linewidth}{0.4pt}
\end{center}
\vspace{14pt}

\begin{longtable}{p{3cm} p{10cm}}
\toprule
\textbf{Abbreviation} & \textbf{Full Form} \\
\midrule
\endhead
ATS & Applicant Tracking System \\
SBERT & Sentence Bidirectional Encoder Representations from Transformers \\
SHAP & SHapley Additive exPlanations \\
O*NET & Occupational Information Network (U.S. Department of Labor) \\
XGBoost & Extreme Gradient Boosting \\
NLP & Natural Language Processing \\
JD & Job Description \\
YOE & Years of Experience \\
CV & Cross-Validation \\
F1 & F1 Score (harmonic mean of Precision and Recall) \\
TF-IDF & Term Frequency -- Inverse Document Frequency \\
DTU & Delhi Technological University \\
API & Application Programming Interface \\
BERT & Bidirectional Encoder Representations from Transformers \\
DiGraph & Directed Graph \\
IS-A & Ontological subsumption relationship \\
LR & Logistic Regression \\
$G$ & Skill ontology graph \\
$V$ & Set of ontology nodes (skills) \\
$E$ & Set of ontology edges (relationships) \\
$S_R$ & Skills extracted from resume \\
$S_J$ & Skills required in job description \\
$w(s_r, s_j)$ & Ontology traversal weight \\
$\vec{v}_R$ & SBERT embedding of resume \\
$\vec{v}_J$ & SBERT embedding of job description \\
\bottomrule
\end{longtable}

% ══════════════════════════════════════════════════════════════════════════════
% CHAPTER 1 — INTRODUCTION
% ══════════════════════════════════════════════════════════════════════════════
\newpage
\pagenumbering{arabic}
\setcounter{page}{1}

\chapter{INTRODUCTION}

\section{Overview and Motivation}

The modern job market operates at a scale that exceeds human processing capacity. A single corporate job posting routinely receives several hundred applications within 48 hours, making automated resume screening operationally necessary. The dominant technology for this screening is the \textbf{Applicant Tracking System (ATS)} --- platforms such as Workday, Greenhouse, and their Indian equivalents embedded in Naukri and LinkedIn Recruiter. These systems score candidate profiles against job descriptions using keyword frequency and rule-based matching, and industry surveys consistently show that 75--88\% of applicants are eliminated at this stage before any human reads their application.

This approach suffers from two fundamental failures. First, it is \textit{semantically blind}: a system that scores ``built distributed backend services with Go'' as zero against a JD requiring ``Golang developer with microservices experience'' has not understood anything --- it has merely counted words. Second, it is \textit{completely opaque}: neither candidates nor recruiters can obtain any structured justification for a screening outcome. Emerging regulations such as the EU AI Act classify employment-related AI systems as ``high-risk'' and mandate explainability, human oversight, and documentation of decision logic.

\textbf{ExplainHire} addresses both failures. It is a \textbf{neurosymbolic AI system} that combines a 434-node skill ontology graph (symbolic AI), a fine-tuned SBERT sentence embedding model (neural AI), and an XGBoost gradient-boosted classifier (learned fusion), delivering a full match/no-match prediction with sentence-level, SHAP-attributed explanations for every decision.

\section{Problem Statement}

Given a resume document $R$ (PDF or DOCX) and a job description $J$ (plain text), produce:

\begin{enumerate}
  \item A binary label $\hat{y} \in \{0, 1\}$ with confidence $p \in [0,1]$
  \item An explanation structure $E$ containing: per-skill match evidence (matched skill, ontological distance, resume evidence sentence), ranked missing skill list ordered by counterfactual impact, and SHAP feature attribution across 14 input features
\end{enumerate}

This is treated as supervised binary classification over a 14-dimensional feature space with post-hoc explainability via SHAP and rule-based evidence extraction.

\section{Research Contributions}

\begin{enumerate}
  \item \textbf{Neurosymbolic matching pipeline}: The first resume--job matching system to simultaneously employ symbolic ontology traversal, fine-tuned neural embeddings, and structural features fused by a learned classifier.
  \item \textbf{Domain-specific skill ontology}: A 434-node directed graph derived from O*NET and extended for Indian software engineering hiring, with an 800-entry alias normalization table.
  \item \textbf{Domain-adapted SBERT}: Fine-tuned \texttt{all-MiniLM-L6-v2} on 887 real resume--JD pairs, adapting the model to technical hiring semantics.
  \item \textbf{Structured explainability}: Sentence-level evidence mapping and counterfactual skill gap analysis with ranked learning resources.
  \item \textbf{State-of-the-art results with public validation}: CV F1 = 0.9045, full-data F1 = 0.951, ResumeAtlas F1 = 0.720 --- highest reported among comparable published systems.
\end{enumerate}

\section{Scope and Limitations}

ExplainHire is validated for technical and software engineering roles in English. The ontology covers technology skills relevant to software development, data science, and machine learning. Extension to non-technical domains, multi-lingual resumes, or candidate ranking across a pool are directions for future work.

\section{Report Organization}

Chapter 2 reviews prior work and relevant background. Chapter 3 describes the system architecture in full. Chapter 4 covers dataset construction and model training. Chapter 5 presents experimental results. Chapter 6 discusses design trade-offs and limitations. Chapter 7 concludes.

% ══════════════════════════════════════════════════════════════════════════════
% CHAPTER 2 — BACKGROUND AND RELATED WORK
% ══════════════════════════════════════════════════════════════════════════════
\chapter{BACKGROUND AND RELATED WORK}

\section{The Recruitment Technology Landscape}

Despite the consequences of ATS filtering --- 75--88\% of applicants rejected before human review --- the core technology underlying most ATS products remains primitive. The standard algorithm parses resumes into fields, extracts keywords from job descriptions, scores each resume by keyword overlap count, and ranks candidates by score. This is \textbf{lexically sensitive but semantically blind}: a candidate who writes ``developed machine learning models for fraud detection'' scores zero against a JD requiring ``AI engineer with experience in anomaly detection,'' despite the substantive equivalence.

\section{Prior Work in Resume--Job Matching}

\subsection{Neural Representation Learning}

\textbf{PJFNN} \cite{zhu2018} (Zhu et al., 2018) introduced a hierarchical CNN learning joint representations of resumes and JDs from Baidu click behavior logs, achieving F1 = 0.800 on a proprietary dataset. It requires large implicit feedback data, provides no explanations, and ignores structured skill relationships.

\textbf{Li et al.} \cite{li2020} (2020) applied BERT with multi-head attention for cross-document interaction, reporting accuracy 0.792. While semantically richer than CNN approaches, the model is a black box.

\subsection{Graph and Multi-View Approaches}

\textbf{Bian et al.} \cite{bian2020} (2020) combined GNN representations with BERT using a co-teaching multi-view network, achieving approximately F1 = 0.78. The graph component is co-occurrence-based rather than a structured domain ontology, and no explainability is provided.

\subsection{Sentence Embedding Approaches}

\textbf{conSultantBERT} \cite{lavi2021} (Lavi et al., 2021) fine-tuned Sentence-BERT on the proprietary Randstad recruitment dataset using a Siamese architecture, achieving F1 = 0.749. It is the closest precedent to our SBERT component; key differences are its proprietary data and absence of symbolic or structural signals.

\subsection{Benchmark Work}

\textbf{Heakl et al.} \cite{heakl2024} (2024) introduced the ResumeAtlas dataset (13,389 resumes, 43 categories) with a TF-IDF + XGBoost baseline achieving macro F1 = 0.610. This is the only large-scale public benchmark in the field, and we use it for external validation.

\subsection{Prior Work Comparison}

Table \ref{tab:priorwork_compare} summarizes the key dimensions on which ExplainHire advances over all prior work.

\begin{table}[h]
\centering
\caption{Feature comparison of ExplainHire with prior systems}
\label{tab:priorwork_compare}
\begin{tabular}{lccccc}
\toprule
\textbf{Feature} & \textbf{PJFNN} & \textbf{Li'20} & \textbf{Bian'20} & \textbf{cBERT} & \textbf{ExplainHire} \\
\midrule
Structured ontology & \texttimes & \texttimes & partial & \texttimes & \checkmark \\
Fine-tuned SBERT & \texttimes & \texttimes & \texttimes & \checkmark & \checkmark \\
Structural features & \texttimes & \texttimes & \texttimes & \texttimes & \checkmark \\
Decision explanation & \texttimes & \texttimes & \texttimes & \texttimes & \checkmark \\
Skill gap analysis & \texttimes & \texttimes & \texttimes & \texttimes & \checkmark \\
Public dataset validation & \texttimes & \texttimes & \texttimes & \texttimes & \checkmark \\
\bottomrule
\end{tabular}
\end{table}

\section{Neurosymbolic Artificial Intelligence}

Neurosymbolic AI combines the complementary strengths of symbolic and neural computing. \textbf{Symbolic AI} represents knowledge as explicit structures (rules, graphs, ontologies) --- interpretable by design but brittle to input variation. \textbf{Neural AI} learns representations from data through gradient optimization --- generalizes well across input variations but is opaque. In ExplainHire: the ontology graph provides structured, interpretable skill reasoning (symbolic); SBERT provides semantic language understanding (neural); XGBoost provides learned fusion that discovers the optimal signal weighting from labeled data; and SHAP provides post-hoc explanation of the learned fusion. This ensures every decision has both a statistically learned justification (SHAP values) and a structural justification (ontology traversal paths and evidence sentences).

\section{Explainability Methods}

\subsection{SHAP (SHapley Additive exPlanations)}

SHAP \cite{lundberg2017} grounds feature attribution in cooperative game theory. The Shapley value of feature $i$ for prediction $f(x)$ is:

\begin{equation}
\phi_i = \sum_{S \subseteq F \setminus \{i\}} \frac{|S|!(|F|-|S|-1)!}{|F|!} \left[ f_{S \cup \{i\}}(x_{S \cup \{i\}}) - f_S(x_S) \right]
\end{equation}

SHAP TreeExplainer computes exact Shapley values for tree ensembles including XGBoost in polynomial time.

\subsection{Ontology-Based Explanation}

Independent of SHAP, the skill ontology provides explanation via traversal path: when ``pytorch'' matches the JD requirement ``machine learning'' via a child-to-parent hop, this is simultaneously the computational evidence for the match and a human-intelligible justification. This form of explanation has no equivalent in any purely neural system.

% ══════════════════════════════════════════════════════════════════════════════
% CHAPTER 3 — SYSTEM ARCHITECTURE
% ══════════════════════════════════════════════════════════════════════════════
\chapter{SYSTEM ARCHITECTURE}

\section{Neurosymbolic Architecture Overview}

Figure \ref{fig:architecture} illustrates how three independent matching signals (symbolic, neural, structural) are computed in parallel and fused by XGBoost to produce a final prediction with full explainability.

\begin{figure}[H]
\centering
\figplaceholder[0.88\textwidth]{fig3\_architecture --- Neurosymbolic Architecture Diagram}
\caption{ExplainHire neurosymbolic architecture: three signals fused by learned classifier}
\label{fig:architecture}
\end{figure}

\section{Design Principles}

ExplainHire is guided by four principles: (1) \textbf{Separation of concerns} --- each layer has a single, well-defined responsibility; (2) \textbf{Signal independence} --- the three signals are computed independently before fusion; (3) \textbf{Explainability by design} --- every component contributing to a decision also contributes to its explanation; (4) \textbf{Graceful degradation} --- the system produces meaningful output even when individual signals return edge-case values.

\section{Pipeline Overview}

Figure \ref{fig:pipeline} shows the complete seven-layer pipeline.

\begin{figure}[H]
\centering
\figplaceholder[0.75\textwidth]{fig1\_pipeline --- Seven-Layer Pipeline Diagram}
\caption{ExplainHire seven-layer pipeline architecture}
\label{fig:pipeline}
\end{figure}

\section{L1 --- Input Validation}

The validation layer accepts only PDF and DOCX files, validates file type by magic bytes, enforces a 5 MB size limit, detects and rejects scanned image PDFs (no embedded text), enforces minimum content length, and validates JD non-emptiness. Any failure returns an HTTP error with a user-facing explanation.

\section{L2 --- Parsing}

PDF text is extracted using PyMuPDF; DOCX using python-docx. From the resume, the parser extracts: (a) \textbf{Skills} --- vocabulary-based scan covering all canonical ontology nodes and alias table entries; (b) \textbf{YOE} --- regex patterns matching ``3 years'', ``3+ years'', ``2 to 4 years'' expressions; (c) \textbf{Education Level} --- degree keywords mapped to ordinal ranks (High School=1 through PhD=5); (d) \textbf{Section Coverage} --- presence/absence of Skills, Experience, Projects, and Education sections. The same techniques are applied to the JD.

\section{L3 --- Normalization}

The normalization layer bridges natural language variation in resume text and the controlled ontology vocabulary. The \textbf{alias table} (800+ entries) maps surface forms to canonical node identifiers:

\begin{lstlisting}[language={}]
{ "nodejs": "node.js", "react.js": "react",
  "ML": "machine_learning", "PyTorch": "pytorch" }
\end{lstlisting}

The table was constructed by manual review of all 144 real resumes. After normalization, all skills from resume and JD are canonical ontology node identifiers.

\section{L4 --- Matching Layer}

\subsection{L4a --- Graph Matcher (Symbolic AI)}

\subsubsection{Ontology Structure}

The skill ontology is a directed graph $G = (V, E)$ where each node $v \in V$ is a canonical skill and each edge $(u, v) \in E$ is an IS-A relationship. The graph contains 434 nodes built from the O*NET technology subtree, manually extended for Indian software engineering contexts (Flutter, FastAPI, Golang, Docker, Kubernetes, Redis, etc.).

Figure \ref{fig:ontology} illustrates an example subgraph.

\begin{figure}[H]
\centering
\figplaceholder[0.88\textwidth]{fig4\_ontology --- Skill Ontology Subgraph}
\caption{Skill ontology subgraph showing IS-A relationships and match weights. A resume listing ``pytorch'' receives weight 0.7 against a JD requiring ``machine learning'' via a child-to-parent hop --- keyword matching scores this 0.}
\label{fig:ontology}
\end{figure}

\subsubsection{Traversal Algorithm}

For each required JD skill $s_j$, the matcher finds the best match among all resume skills using the six-level weight scheme:

\begin{equation}
w(s_r, s_j) = \begin{cases}
1.0 & \text{exact match} \\
0.9 & \text{alias match} \\
0.7 & \text{child to parent (resume skill is more specific)} \\
0.5 & \text{parent to child (resume skill is more general)} \\
0.4 & \text{sibling (common ancestor)} \\
0.3 & \text{two-hop traversal} \\
0.0 & \text{no connection}
\end{cases}
\end{equation}

The overall graph score is:
\begin{equation}
\text{GraphScore}(R, J) = \frac{1}{|S_J|} \sum_{s_j \in S_J} \max_{s_r \in S_R} w(s_r, s_j)
\end{equation}

The graph matcher contributes 5 features to the feature vector: \{graph\_score, coverage, exact\_matches, partial\_matches, no\_matches\}.

\subsection{L4b --- Semantic Matcher (Neural AI)}

The semantic matcher employs \texttt{all-MiniLM-L6-v2} SBERT \cite{reimers2019}, a 6-layer distilled Siamese BERT with 384-dimensional embeddings. The semantic similarity is:

\begin{equation}
\text{SBERT\_Score}(R, J) = \frac{\vec{v}_R \cdot \vec{v}_J}{\|\vec{v}_R\| \cdot \|\vec{v}_J\|}
\end{equation}

Long texts are truncated to the last 2,000 characters before encoding, prioritizing skills/projects sections over boilerplate headers.

\textbf{Fine-Tuning}: The base model was fine-tuned using CosineSimilarityLoss on 887 real resume--JD pairs:
\begin{equation}
\mathcal{L} = \frac{1}{N} \sum_{i=1}^{N} \left( \cos(\vec{v}_{R_i}, \vec{v}_{J_i}) - y_i \right)^2
\end{equation}

where $y_i \in \{0.0, 1.0\}$ is the match label. Training ran for 4 epochs with batch size 16, reaching final loss 0.069. This teaches the model what ``similar'' means in a hiring context rather than general English.

\subsection{L4c --- Structural Matcher}

\textbf{YOE Score:}
\begin{equation}
\text{YOE\_Score} = \max\!\left(0,\ 1 - \frac{\max(0,\; \text{YOE}_\text{req} - \text{YOE}_\text{resume})}{\text{YOE}_\text{req} + 1}\right)
\end{equation}

\textbf{Education Score:} Education levels are ranked 1--5. Score is 1.0 if candidate meets or exceeds requirement, 0.5 if one rank below, 0.0 otherwise.

\textbf{Section Score:}
\begin{equation}
\text{Section\_Score} = \frac{|\text{detected sections}|}{4}
\end{equation}

The structural matcher contributes 8 features: \{structural\_score, yoe\_score, edu\_score, section\_score, resume\_yoe, required\_yoe, resume\_edu\_rank, required\_edu\_rank\}.

\section{L5 --- Classification and Explainability}

\subsection{Feature Vector}

The 14-dimensional feature vector fused by XGBoost is:
\begin{equation}
\mathbf{x} = [\underbrace{g_1,\ldots,g_5}_{\text{graph}},\ \underbrace{s_1}_{\text{SBERT}},\ \underbrace{t_1,\ldots,t_8}_{\text{structural}}]
\end{equation}

\subsection{XGBoost Classifier}

XGBoost \cite{chen2016} builds an additive ensemble of decision trees where each tree corrects the residual errors of its predecessors:
\begin{equation}
\hat{y}^{(t)} = \hat{y}^{(t-1)} + \eta \cdot f_t(\mathbf{x})
\end{equation}

XGBoost was selected for three reasons: state-of-the-art performance on small tabular datasets; native SHAP-compatible tree structures; and CPU-only training feasible in seconds.

\subsection{SHAP Attribution}

After XGBoost produces its prediction, SHAP TreeExplainer computes exact Shapley values $\phi_i$ for each of the 14 features:
\begin{equation}
f(\mathbf{x}) = \phi_0 + \sum_{i=1}^{14} \phi_i
\end{equation}

Figure \ref{fig:shap} shows a representative SHAP attribution plot. Positive $\phi_i$ pushes toward Match; negative pushes toward No Match.

\begin{figure}[H]
\centering
\figplaceholder[0.85\textwidth]{fig5\_shap --- SHAP Feature Attribution Bar Chart}
\caption{SHAP feature attribution for an example prediction. Strong graph and SBERT scores drive the Match verdict despite a YOE shortfall.}
\label{fig:shap}
\end{figure}

\section{L6 --- Explainer}

\textbf{Evidence Mapper}: For each matched JD skill, the mapper finds the most relevant resume sentence via TF-IDF weighted token overlap. The user sees the matched skill, its ontological distance, and the specific resume sentence as evidence.

\textbf{Skill Gap Ranker}: For each unmatched skill, the ranker estimates counterfactual impact:
\begin{equation}
\Delta_j = P_\text{XGB}(\mathbf{x}'_j) - P_\text{XGB}(\mathbf{x})
\end{equation}

where $\mathbf{x}'_j$ is the feature vector with the missing skill hypothetically added. Gaps are ranked by $\Delta_j$ and presented with curated learning resource links.

\section{L7 --- Web Application}

The Flask application provides: resume upload (PDF/DOCX), JD text input, a results dashboard with match verdict, three score cards, per-skill breakdown with evidence sentences, ranked skill gap list, and SHAP bar chart. Additional features include match history (SQLite/Flask-SQLAlchemy), user authentication (Flask-Login, bcrypt), an interactive 434-node ontology visualizer (vis-network), and dark/light mode toggle.

% ══════════════════════════════════════════════════════════════════════════════
% CHAPTER 4 — DATASET AND TRAINING
% ══════════════════════════════════════════════════════════════════════════════
\chapter{DATASET AND TRAINING}

\section{Data Collection Strategy}

Labeled resume--JD pairs are not publicly available at scale. We assembled a multi-source dataset combining synthetic edge cases, a public Kaggle collection, and a proprietary dataset of real resumes from our peer network matched against real industry JDs.

\section{Dataset Sources}

\subsection{Synthetic Pairs (167 rows)}

Manually constructed to cover: semantic equivalence pairs (testing SBERT), near-miss negative pairs (testing false positive resistance), ontology boundary cases (testing sibling/child matching), and experience gap pairs (testing structural component). All pairs were independently labeled by two team members with disagreements resolved by discussion.

\subsection{Kaggle Resume Dataset (800 rows)}

The Kaggle Resume Dataset (Anbhawal, 2020) contains 800 real resumes across 24 job categories. We constructed resume--JD pairs by matching resumes against JDs from their own category (positive) and randomly sampled different categories (negative).

\subsection{Friends Dataset (887 rows)}

We collected 144 real B.Tech student resumes and matched them against 17 real industry JDs from Naukri and LinkedIn, covering roles including Full Stack Developer, Data Scientist, Backend Engineer, Flutter Developer, Machine Learning Engineer, Golang Developer, DevOps Engineer, and Android Developer. The full 144 $\times$ 17 = 2,448 cross-product was feature-extracted and score-based auto-labeled.

\begin{table}[h]
\centering
\caption{Training dataset composition}
\label{tab:dataset}
\begin{tabular}{lrrl}
\toprule
\textbf{Source} & \textbf{Rows} & \textbf{Label Type} & \textbf{Notes} \\
\midrule
Synthetic pairs & 167 & Manual annotation & Edge cases \\
Kaggle Resume Dataset & 800 & Category-based & 24 job categories \\
Friends dataset & 887 & Score-based auto-label & 144 resumes $\times$ 17 JDs \\
\midrule
\textbf{Total} & \textbf{1,854} & & \\
\bottomrule
\end{tabular}
\end{table}

\section{Labeling Methodology}

For the friends dataset, manual labeling at 2,448 pairs was infeasible. We applied conservative threshold-based auto-labeling requiring both signals to agree strongly:

\begin{equation}
\text{label}(R, J) = \begin{cases}
1 & \text{if } \text{GraphScore} > 0.4 \text{ AND } \text{SBERT\_Score} > 0.4 \\
0 & \text{if } \text{GraphScore} < 0.3 \text{ AND } \text{SBERT\_Score} < 0.3 \\
-1 & \text{excluded (signals disagree)}
\end{cases}
\end{equation}

This yielded 636 positive and 251 negative pairs (887 total), with 1,561 mixed pairs excluded. The label distribution is 71.7\% positive / 28.3\% negative.

\section{SBERT Fine-Tuning}

Fine-tuning used the sentence-transformers library with CosineSimilarityLoss, batch size 16, 4 epochs, 88 warmup steps, and AdamW optimizer. Final training loss: \textbf{0.069}. The fine-tuned model was saved to \texttt{models/sbert\_finetuned/} and is automatically loaded by the pipeline.

\begin{table}[h]
\centering
\caption{Effect of SBERT fine-tuning on classifier performance}
\begin{tabular}{lcc}
\toprule
\textbf{Configuration} & \textbf{CV F1} & \textbf{Full-data F1} \\
\midrule
Base SBERT + 967 training rows & 0.899 & 0.899 \\
Fine-tuned SBERT + 1,854 rows & \textbf{0.9045} & \textbf{0.951} \\
\bottomrule
\end{tabular}
\end{table}

\section{XGBoost Training}

The 14-feature vector was extracted for all 1,854 pairs and cached in \texttt{data/annotation.csv}. XGBoost was trained with the following configuration:

\begin{table}[h]
\centering
\caption{XGBoost hyperparameters}
\begin{tabular}{ll}
\toprule
\textbf{Parameter} & \textbf{Value} \\
\midrule
n\_estimators & 200 \\
max\_depth & 6 \\
learning\_rate & 0.1 \\
tree\_method & hist (CPU-optimized) \\
eval\_metric & mlogloss \\
random\_state & 42 \\
\bottomrule
\end{tabular}
\end{table}

Model performance was evaluated using 5-fold stratified cross-validation (stratified to preserve the 71.7/28.3\% class ratio in each fold). The CV F1 of \textbf{0.9045} is our primary honest performance estimate.

% ══════════════════════════════════════════════════════════════════════════════
% CHAPTER 5 — EXPERIMENTS AND RESULTS
% ══════════════════════════════════════════════════════════════════════════════
\chapter{EXPERIMENTS AND RESULTS}

\section{Experimental Setup}

All experiments ran on a consumer laptop (Intel Core i5, 16 GB RAM, no GPU), demonstrating that ExplainHire is deployable without specialized hardware. SBERT fine-tuning completed in approximately 15 minutes on CPU. All F1 scores are macro-averaged using scikit-learn's \texttt{f1\_score(average="macro", zero\_division=0)}, which penalizes poor performance on the minority class.

\section{Ablation Study}

We evaluated five configurations on all 1,854 labeled pairs using 5-fold stratified cross-validation.

\begin{figure}[H]
\centering
\figplaceholder[0.90\textwidth]{fig2\_ablation --- Ablation Study Bar Chart}
\caption{Ablation study: macro F1 by configuration. Full ExplainHire pipeline achieves F1 = 0.951, improving +11.86\% over the best single-signal baseline.}
\label{fig:ablation}
\end{figure}

\begin{table}[h]
\centering
\caption{Ablation study results (5-fold stratified CV, 1,854 pairs)}
\label{tab:ablation}
\begin{tabular}{lccc}
\toprule
\textbf{Configuration} & \textbf{Macro F1} & \textbf{$\Delta$ vs. prev.} & \textbf{Interpretation} \\
\midrule
TF-IDF baseline & 0.037 & --- & Keyword matching fails entirely \\
SBERT-only & 0.484 & +44.7\% & Semantics help; global similarity imprecise \\
Graph-only & 0.747 & +26.3\% & Ontology is the strongest single signal \\
Graph + SBERT & 0.832 & +8.5\% & Symbolic + neural are complementary \\
\textbf{Full pipeline} & \textbf{0.951} & \textbf{+11.9\%} & Structural features add independent value \\
\bottomrule
\end{tabular}
\end{table}

The TF-IDF baseline (F1 = 0.037) confirms keyword co-occurrence is nearly useless. The Graph-only result (F1 = 0.747) is the most informative single result: structured skill matching over an ontology is substantially more informative than either keyword matching or global semantic similarity, validating the core architectural decision. The full pipeline (F1 = 0.951) confirms all three signals contribute independently.

\section{External Benchmark: ResumeAtlas}

To assess generalizability, we evaluated on the public ResumeAtlas dataset (13,389 resumes, 43 job categories). We trained on 80\% and evaluated on 20\% (2,678 resumes).

\begin{table}[h]
\centering
\caption{ResumeAtlas benchmark results}
\label{tab:benchmark}
\begin{tabular}{lcc}
\toprule
\textbf{Method} & \textbf{Macro F1} & \textbf{vs. Heakl et al. baseline} \\
\midrule
TF-IDF + LR \textit{(Heakl et al.)} & 0.610 & --- \\
SBERT + LR & 0.706 & +9.6\% \\
\textbf{SBERT + XGBoost (ExplainHire)} & \textbf{0.720} & \textbf{+11.06\%} \\
\bottomrule
\end{tabular}
\end{table}

\section{Comparison with Prior Work}

\begin{table}[h]
\centering
\caption{Comparison with published prior work}
\label{tab:priorwork}
\begin{tabular}{llccc}
\toprule
\textbf{System} & \textbf{Method} & \textbf{F1} & \textbf{Dataset} & \textbf{Explainable} \\
\midrule
PJFNN \cite{zhu2018} & CNN joint representation & 0.800 & Baidu (private) & No \\
Li et al. \cite{li2020} & BERT + attention & 0.792 & CRC (private) & No \\
Bian et al. \cite{bian2020} & Graph + BERT multi-view & $\sim$0.78 & Private & No \\
conSultantBERT \cite{lavi2021} & Fine-tuned SBERT & 0.749 & Randstad (private) & No \\
Heakl et al. \cite{heakl2024} & TF-IDF + XGBoost & 0.610 & ResumeAtlas & No \\
\midrule
\textbf{ExplainHire (full)} & \textbf{Neurosymbolic} & \textbf{0.951} & \textbf{This work} & \textbf{Yes} \\
ExplainHire on ResumeAtlas & SBERT + XGBoost & 0.720 & ResumeAtlas & Yes \\
\bottomrule
\end{tabular}
\end{table}

ExplainHire achieves the highest reported F1 among all comparable systems, and is the only system providing structured, auditable explanations.

% ══════════════════════════════════════════════════════════════════════════════
% CHAPTER 6 — DISCUSSION
% ══════════════════════════════════════════════════════════════════════════════
\chapter{DISCUSSION}

\section{Signal Importance}

The ablation results reveal a hierarchy: Graph (0.747) $>$ SBERT (0.484) $>$ Structural (+11.9\% at margin). The graph is the most valuable single component because resume--job matching is fundamentally a structured coverage problem: a job requires specific skills, a resume has specific skills, and the question is whether the candidate's skills cover the requirements. A structured knowledge graph is the natural representation for hierarchical coverage checking, and the 6-level traversal credits candidates with related skills as a human reviewer would. SBERT captures contextual alignment the graph cannot, but global similarity is an imprecise proxy for skill-by-skill coverage. Structural features add orthogonal information --- a candidate can have all required skills but insufficient experience, which neither ontology nor SBERT detects.

\section{Comparison with Large Language Models}

LLMs (GPT-4, Claude, Gemini) can assess resume--JD fit through prompting. ExplainHire's advantages over LLMs in this context are: (1) \textbf{Consistency} --- same input always produces same output, essential for fair employment screening; (2) \textbf{Cost and speed} --- milliseconds at near-zero marginal cost versus expensive per-token LLM inference at scale; (3) \textbf{Structured explainability} --- SHAP values and ontology traversal paths constitute an audit trail, unlike LLM natural language reasoning; (4) \textbf{Regulatory compliance} --- feature-level attribution is auditable under emerging employment AI regulations. ExplainHire is not competing with LLMs for nuanced text understanding --- it is competing with ATS systems for structured, explainable, scalable employment screening.

\section{Limitations}

\begin{itemize}[noitemsep]
  \item \textbf{Dataset scope}: Training data biased toward Indian B.Tech students in software engineering. Performance on non-technical domains or senior roles is unevaluated.
  \item \textbf{Auto-labeling noise}: Score-based thresholding of friends-dataset pairs may introduce label errors despite conservative thresholds.
  \item \textbf{Ontology gaps}: Any JD skill absent from the 434-node graph will receive graph score 0 regardless of semantic equivalence in the resume.
  \item \textbf{YOE extraction}: Regex-based extraction may underestimate YOE for candidates stating experience implicitly rather than explicitly.
\end{itemize}

\section{Future Work}

Directions include: ontology expansion to non-technical domains; active learning using counterfactual gap ranking to prioritize manual labeling; multi-lingual resume support; multi-candidate ranking; and LLM integration as a fourth signal for narrative understanding.

% ══════════════════════════════════════════════════════════════════════════════
% CHAPTER 7 — CONCLUSION
% ══════════════════════════════════════════════════════════════════════════════
\chapter{CONCLUSION}

This report presented \textbf{ExplainHire}, a neurosymbolic AI system for explainable resume--job description matching. The seven-layer pipeline combines a 434-node symbolic skill ontology graph, a domain-adapted SBERT model fine-tuned on 887 real recruitment pairs, and structural features (YOE, education, section coverage), fused by XGBoost with SHAP explainability and sentence-level evidence mapping.

Key findings:
\begin{enumerate}
  \item All three signals are necessary. No single signal approaches full-pipeline performance (TF-IDF: 0.037; full pipeline: 0.951).
  \item The symbolic ontology graph is the most valuable single component (F1 = 0.747), outperforming neural SBERT (F1 = 0.484) and confirming that structured domain knowledge cannot be replaced by statistical language models for skill matching.
  \item Structural features add +11.9\% F1 at the margin, demonstrating significant independent predictive power beyond the two skill-matching signals.
  \item ExplainHire generalizes to public benchmarks: F1 = 0.720 on ResumeAtlas vs. 0.610 reported baseline (+11.06\%).
  \item Explainability is achieved without sacrificing accuracy. ExplainHire is simultaneously the most accurate and the only explainable system among all comparable published works.
\end{enumerate}

ExplainHire demonstrates that the neurosymbolic paradigm is a productive approach for high-stakes NLP tasks requiring both accuracy and accountability. The recruitment domain is one of many such applications; similar architectures could extend to medical diagnosis, legal document review, and other domains where decisions must be both accurate and auditable.

% ══════════════════════════════════════════════════════════════════════════════
% APPENDICES
% ══════════════════════════════════════════════════════════════════════════════
\appendix

\chapter{Project File Structure}

\begin{lstlisting}[language={}]
ExplainHire/
|-- app/                         # Flask web application
|   |-- __init__.py
|   |-- models.py                # SQLAlchemy DB models
|   `-- routes/
|       |-- auth.py              # /login /register /logout
|       |-- match.py             # /match  (main pipeline)
|       |-- history.py           # /history /result/<id>
|       `-- ontology_view.py     # /ontology
|-- pipeline/
|   |-- l1_input/
|   |-- l2_parse/
|   |-- l3_normalize/
|   |-- l4_match/
|   |   |-- graph_matcher.py
|   |   |-- semantic_matcher.py
|   |   |-- structural_matcher.py
|   |   `-- finetune_sbert.py
|   |-- l5_classify/
|   |   |-- trainer.py
|   |   `-- predictor.py
|   |-- l6_explain/
|   |   |-- evidence_mapper.py
|   |   `-- skill_gap.py
|   `-- l7_present/
|       `-- orchestrator.py
|-- ontology/
|   |-- skill_ontology.gpickle   # 434-node NetworkX graph
|   |-- alias_table.json         # 800+ surface forms
|   `-- learning_resources.json
|-- data/
|   |-- raw/friends/             # 144 real B.Tech resumes
|   |-- raw/jds/                 # 17 real industry JDs
|   |-- annotation.csv           # 1,854 labeled pairs
|   `-- friends_rows.csv
|-- evaluation/
|   |-- ablation.py
|   |-- baselines.py
|   `-- benchmark_resumeatlas.py
|-- models/
|   |-- xgb_matcher.pkl
|   `-- sbert_finetuned/
|-- paper/
|-- config.py
|-- requirements.txt
`-- run.py
\end{lstlisting}

\chapter{Key Numerical Results}

\begin{table}[h]
\centering
\caption{Summary of all key numerical results}
\begin{tabular}{ll}
\toprule
\textbf{Metric} & \textbf{Value} \\
\midrule
Pipeline layers & 7 \\
Ontology nodes & 434 \\
Alias table entries & 800+ \\
Total training pairs & 1,854 \\
SBERT fine-tuning pairs & 887 \\
SBERT fine-tuning loss & 0.069 \\
SBERT fine-tuning epochs & 4 \\
XGBoost trees & 200 \\
Feature vector size & 14 \\
Cross-validated F1 & 0.9045 \\
Full-dataset F1 & 0.951 \\
Best single-signal F1 (Graph-only) & 0.747 \\
Improvement over best baseline & +11.86\% \\
TF-IDF baseline F1 & 0.037 \\
ResumeAtlas F1 & 0.720 \\
Improvement over Heakl et al. & +11.06\% \\
\bottomrule
\end{tabular}
\end{table}

% ══════════════════════════════════════════════════════════════════════════════
% BIBLIOGRAPHY
% ══════════════════════════════════════════════════════════════════════════════
\newpage
\begin{thebibliography}{99}

\bibitem{zhu2018}
R. Zhu et al., ``Person-Job Fit: Adapting the Right Talent for the Right Job with Joint Representation Learning,'' \textit{ACM TMIS}, vol.~9, no.~3, 2018. DOI: \texttt{10.1145/3234465}

\bibitem{li2020}
X. Li et al., ``BERT-based Neural Collaborative Filtering and Fixed-point Number Embedding for Job-Resume Matching,'' \textit{EMNLP 2020}. DOI: \texttt{10.18653/v1/2020.emnlp-main.681}

\bibitem{bian2020}
S. Bian et al., ``Learning to Match Jobs with Resumes from Sparse Interaction Data using Multi-View Co-Teaching Network,'' \textit{CIKM 2020}. DOI: \texttt{10.1145/3340531.3411929}

\bibitem{lavi2021}
D. Lavi, O. Medini, and O. Kurland, ``consultantBERT: Fine-tuned Siamese Sentence-BERT for Job Title Benchmarking,'' \textit{arXiv:2109.06912}, 2021.

\bibitem{heakl2024}
A. Heakl et al., ``ResumeAtlas: Revisiting Resume Classification with Large-Scale Datasets and Large Language Models,'' \textit{arXiv:2406.18125}, 2024.

\bibitem{reimers2019}
N. Reimers and I. Gurevych, ``Sentence-BERT: Sentence Embeddings using Siamese BERT-Networks,'' \textit{EMNLP 2019}. DOI: \texttt{10.18653/v1/D19-1410}

\bibitem{chen2016}
T. Chen and C. Guestrin, ``XGBoost: A Scalable Tree Boosting System,'' \textit{KDD 2016}. DOI: \texttt{10.1145/2939672.2939785}

\bibitem{lundberg2017}
S. Lundberg and S.-I. Lee, ``A Unified Approach to Interpreting Model Predictions,'' \textit{NeurIPS 2017}. DOI: \texttt{10.48550/arXiv.1705.07874}

\bibitem{onet}
National Center for O*NET Development, ``O*NET OnLine,'' U.S. Department of Labor. [Online]. Available: \url{https://www.onetonline.org}

\end{thebibliography}

% ── Proof pages ───────────────────────────────────────────────────────────────
\newpage
\thispagestyle{plain}
\begin{center}
\vspace*{5cm}
{\Large \textbf{PAPER ACCEPTANCE PROOF}}\\[20pt]
\textit{(To be attached upon paper acceptance)}
\end{center}

\newpage
\thispagestyle{plain}
\begin{center}
\vspace*{5cm}
{\Large \textbf{REGISTRATION PROOF}}\\[20pt]
\textit{(To be attached upon conference registration)}
\end{center}

\newpage
\thispagestyle{plain}
\begin{center}
\vspace*{5cm}
{\Large \textbf{SCOPUS INDEXED CONFERENCE PROOF}}\\[20pt]
\textit{(To be attached)}
\end{center}

\end{document}
